% Generated draft: requires mathematical and source review.
\documentclass[11pt]{article}
\usepackage[margin=25mm]{geometry}
\usepackage{fontspec}
\setmainfont{Latin Modern Roman}
\usepackage{amsmath,amssymb,mathtools,unicode-math}
\setmathfont{Latin Modern Math}
\usepackage{enumitem,needspace,xurl,hyperref}
\setlength{\parindent}{0pt}
\setlength{\parskip}{5pt}
\setlength{\emergencystretch}{3em}
\newcommand{\N}{\mathbb{N}}
\newcommand{\Z}{\mathbb{Z}}
\newcommand{\Q}{\mathbb{Q}}
\newcommand{\R}{\mathbb{R}}
\newcommand{\C}{\mathbb{C}}
\newcommand{\F}{\mathbb{F}}
\newcommand{\A}{\mathbb{A}}
\newcommand{\PP}{\mathbb{P}}
\newcommand{\E}{\mathbb{E}}
\newcommand{\eps}{\varepsilon}
\newcommand{\dd}{\,\mathrm{d}}
\begin{document}
\section*{UnsolvedMath ID: 800376}
\textbf{Source problem code:} catalogue entry
\subsection{Definitions and mathematical statement}
Write \(\mathbb{N}_0=\{0,1,2,\ldots\}\) and \(\mathbb{P}=\{1,2,3,\ldots\}\). We interpret the source's indexing set \(\mathbb{N}\) as \(\mathbb{N}_0\); the supplied excerpt does not explicitly state this convention. For \(m\in\mathbb{P}\), let \(S_m\) be the set of permutations of \(\{1,\ldots,m\}\). For \(\pi\in S_m\), define
\[
\operatorname{des}(\pi)=\left|\{i\in\{1,\ldots,m-1\}:\pi(i)>\pi(i+1)\}\right|.
\]
For \(k\in\mathbb{N}_0\), the Eulerian number is
\[
A(m,k)=\left|\{\pi\in S_m:\operatorname{des}(\pi)=k\}\right|.
\]
Thus \(A(m,k)=0\) for \(k\ge m\), and \(A(1,0)=1\). Define the infinite matrix
\[
A=\bigl(A(n+1,k)\bigr)_{n,k\in\mathbb{N}_0}.
\]
Here the source's term \emph{totally positive} is interpreted as requiring every finite minor to be nonnegative, a property also called total nonnegativity. The excerpt does not define this terminology, so this interpretation requires confirmation; requiring every minor to be strictly positive would be incompatible with the matrix's zero entries.

\textbf{Conjecture (Brenti, source Conjecture 14).} With these indexing and terminology conventions, for every integer \(r\ge1\), every sequence \(0\le n_1<\cdots<n_r\), and every sequence \(0\le k_1<\cdots<k_r\),
\[
\det\!\bigl(A(n_i+1,k_j)\bigr)_{1\le i,j\le r}\ge0.
\]
The supplied source attributes the first formulation of this conjecture to [BreCombTP, Conjecture 6.10].
\subsection{Short English statement}
Brenti's supplied conjecture asks whether every finite minor of the infinite Eulerian-number matrix is nonnegative. This formulation interprets the source's natural-number indices as starting at zero and its term ``totally positive'' in the nonnegative-minor sense; both conventions require source confirmation.
\subsection{Sources}
\begin{itemize}
\item Local UnsolvedMath registry, permanent ID \(800376\), exact source problem code \texttt{catalogue entry}, title ``Coxeter Groups and Unimodality --- Conjecture 14.'' The supplied record attributes the problem to Brenti, with proposed year 2024, and labels it open. This status is supplied registry information, not an independently verified literature assessment.
\item Brenti, \emph{Some Open Problems on Coxeter Groups and Unimodality} (2024), Conjecture 14. Supplied source URL: \url{https://arxiv.org/abs/2410.09897}. The supplied provenance records access on 2026-10-16 and extraction by \texttt{arxiv\_latex}; no independent source visit is claimed here.
\item [BreCombTP, Conjecture 6.10], identified in the supplied excerpt as the first formulation of the conjecture. The record does not supply the full bibliographic entry or a separate URL.
\item [BreMem], [BreCombTP], and [BreTPComb], cited in the supplied excerpt for background on totally positive matrices in combinatorics. Their full bibliographic entries are not supplied.
\end{itemize}
\end{document}
